\chapter{因果发现}
\label{chapter::causal-discovery}

\section{引言：一个雄心勃勃的问题}

在前面的章节中，我们假设因果结构已知——即已知哪些变量是原因、哪些变量是结果。然而在现实中，我们往往面对一个更为艰难的问题：\textbf{能否从数据中学习出变量之间的因果关系？}

这个问题的重要性不言而喻。科学发现的本质就是从观察数据中推断因果结构。历史上，James Lind 通过对照实验发现柑橘类水果可以治疗坏血病；但在许多情况下，实验干预是不可行或不道德的。我们能否仅从被动观察的数据中，发现变量之间的因果关系？

这个问题在哲学上曾被认为是不可能的。``相关不蕴含因果''——这句格言道出了统计学的传统智慧。然而，过去三十年见证了一场静默的革命：在一系列精心设计的假设下，因果结构可以从纯粹的观察数据中学习。

这场革命的先驱包括 Judea Pearl，他发展了基于图的因果推理理论；以及 Peter Spirtes、Clark Glymour 和 Richard Scheines，他们构建了 PC 等算法从数据中提取因果结构。本章将介绍这一领域的基本思想和方法。

\section{因果马尔可夫条件与忠诚性}

在讨论具体算法之前，我们需要理解因果发现的可能性条件。考虑一个包含变量 $X_1, X_2, \ldots, X_p$ 的系统。如果这些变量之间存在因果关系，我们通常用有向无环图（DAG）来表示：每个节点代表一个变量，箭头 $X \to Y$ 表示 $X$ 是 $Y$ 的直接原因（或之一）。

一个根本性的观察是：在许多合理的因果模型下，变量之间的条件独立性蕴含着因果结构的信息。

\begin{assumption}[因果马尔可夫条件]
\label{assume::causal-markov}
给定一个变量的所有直接原因，该变量独立于所有非其后代的变量。
\end{assumption}

直观上，如果 $X$ 导致 $Y$，而 $Y$ 又导致 $Z$，那么在已知 $Y$ 的情况下，$X$ 和 $Z$ 之间不应该有直接的联系——它们之间的任何关联都已经被 $Y$ 解释了。

\begin{assumption}[忠诚性条件]
\label{assume::faithfulness}
图中存在的条件独立性关系，恰好对应于数据分布中存在的条件独立性关系。
\end{assumption}

忠诚性条件防止了``巧合独立性''——即两个变量在数据中条件独立，但在图中却通过一条路径相连。这种巧合在参数化模型中是零测度的。

在这两个假设下，我们可以从变量的条件独立性结构中读取因果信息。

\section{PC 算法：从条件独立性到因果骨架}

PC 算法（Spirtes et al., 2000）是因果发现领域最经典的算法之一。它的核心思想简洁而有力：\textbf{从条件独立性出发，逐步确定因果结构。}

\subsection{第一步：发现骨架}

算法从一个完全无向的图出发，通过条件独立性检验逐步删除边。

对于每一对变量 $X$ 和 $Y$，以及越来越大的条件集 $S$（最初为空，然后是单变量，然后是双变量，以此类推），我们检验条件独立性：
\begin{equation}
\label{eq:cond-indep-test}
X \Perp Y \mid S
\end{equation}
其中 $\Perp$ 表示条件独立性。如果拒绝独立性，则保留边 $(X, Y)$；否则，删除这条边。

在标准的高斯或离散数据设置下，我们可以使用偏相关系数或条件互信息来进行这个检验。

\begin{Example}[简单的条件独立性]
考虑三个变量 $X, Y, Z$。假设我们进行条件独立性检验：
\begin{itemize}
\item $X \not\Perp Y$（边缘不独立）
\item $X \Perp Y \mid Z$（给定 $Z$ 条件独立）
\end{itemize}
这告诉我们什么？在 Markov 条件下，这暗示 $Z$ 是 $X$ 和 $Y$ 的共同原因，或者 $Z$ 位于它们之间的中介路径上。
\end{Example}

这个过程不断迭代，直到所有边都经过了适当的条件独立性检验。最终得到的图称为\textbf{因果骨架}——它保留了所有在某种条件下不条件独立的变量对之间的连接。

\subsection{第二步：定向边}

骨架是无向的，但我们需要定向边来揭示因果方向。这需要利用 v-结构（碰撞结构）的概念。

考虑三个变量 $X, Y, Z$，它们之间的骨架显示 $X-Z-Y$（即 $X$ 与 $Y$ 之间无边，但都与 $Z$ 相连）。如果进一步发现 $X \not\Perp Y \mid \emptyset$ 且 $X \Perp Y \mid Z$，这符合什么结构？

在 Markov 条件下，如果 $X \to Z \to Y$（$Z$ 是中介），那么给定 $Z$，$X$ 和 $Y$ 应该条件独立。这与观察到的 $X \Perp Y \mid Z$ 一致。

但如果有 $Z$ 是 $X$ 和 $Y$ 的共同原因，即 $Z \to X$ 且 $Z \to Y$，那么 $X$ 和 $Y$ 在边缘上不独立，但给定 $Z$ 时条件独立——这正是我们观察到的！

PC 算法通过识别这类模式来确定边的方向。形式化地说：

\begin{definition}[v-结构]
\label{def:v-structure}
一个三元组 $(X, Z, Y)$ 形成 v-结构，当且仅当：
\begin{enumerate}
\item $X$ 与 $Y$ 在骨架中无边；
\item $Z$ 与 $X$、$Z$ 与 $Y$ 都有边相连；
\item 在给定 $Z$ 的条件下，$X$ 和 $Y$ 不条件独立。
\end{enumerate}
\end{definition}

在忠诚性条件下，v-结构必然对应一个碰撞结构：$X \to Z \gets Y$。这是因为其他结构（$X \to Z \to Y$ 或 $X \gets Z \to Y$）都会在给定 $Z$ 时使 $X$ 和 $Y$ 条件独立。

PC 算法的定向规则可以总结为：
\begin{enumerate}
\item 识别所有 v-结构，将其中心节点的两个父节点方向设为指向该节点；
\item 应用马尔可夫性和忠诚性所隐含的其他定向规则（如避免产生新的 v-结构）。
\end{enumerate}

\subsection{第三步：边定向的最终确定}

上述步骤可能无法定向所有边。这是因为某些边在观察分布下本来就是不可分辨的——不同的有向图可能产生完全相同的条件独立性结构。这些图构成一个\textbf{等价类}。

例如，考虑两条边 $X - Y$ 和 $Y - Z$，且 $X$ 与 $Z$ 之间无边。$X \to Y \to Z$ 和 $X \gets Y \gets Z$ 产生完全相同的条件独立性结构，在没有额外假设的情况下无法区分。

这引导我们到下一个重要概念。

\section{DAG 等价类与基本等价类图}

\textbf{部分有向无环图（PAG）}是一种表示 DAG 等价类的标准形式。在 PAG 中：
\begin{itemize}
\item 不可定向的边用 $\circ - \circ$ 表示；
\item 可定向的边用 $\to$ 或 $\gets$ 表示。
\end{itemize}

\begin{Theorem}[忠实性定理的推论]
\label{eq:pag-uniqueness}
在因果马尔可夫条件和忠诚性条件下，观察分布唯一确定一个 PAG。
\end{Theorem}

这个定理保证了 PC 算法输出的 PAG 是唯一的——如果我们相信这些假设的话。

\section{条件独立性检验}

PC 算法的核心依赖于可靠的条件独立性检验。这在实践中是一个 nontrivial 的问题。

\subsection{高斯情形：偏相关系数}

对于连续型高斯数据，条件独立性可以通过偏相关系数来检验。给定条件集 $S$，$X$ 和 $Y$ 的偏相关系数记为 $\rho_{XY \cdot S}$。

我们可以检验零假设 $H_0: \rho_{XY \cdot S} = 0$。在正态性假设下，这个统计量服从一个已知分布，可以计算 p 值。

\begin{Remark}
偏相关系数的计算需要 $S$ 中变量的条件协方差矩阵。当 $|S|$ 较大时，估计精度会下降，这是 PC 算法的一个 practical limitation。
\end{Remark}

\subsection{离散情形：条件互信息}

对于离散数据，我们可以使用条件互信息：
\begin{equation}
\label{eq:cond-mutual-info}
I(X; Y \mid S) = \sum_{x, y, s} p(x, y, s) \log \frac{p(x, y \mid s)}{p(x \mid s) p(y \mid s)}.
\end{equation}
在独立性下，这个量应该为零。我们可以使用卡方检验或似然比检验来验证。

\section{PC 算法的完整流程}

综合以上讨论，PC 算法的完整流程如下：

\begin{enumerate}
\item \textbf{初始化}：从一个完全无向的完全图开始。
\item \textbf{骨架发现}：对每一对变量 $(X, Y)$，以及越来越大的条件集 $S$，进行条件独立性检验。如果 $X \Perp Y \mid S$，则删除边 $(X, Y)$。
\item \textbf{v-结构识别}：对于每个三元组 $(X, Z, Y)$，检查是否形成 v-结构。如果是，定向 $X \to Z \gets Y$。
\item \textbf{边定向传播}：应用定向规则，尽可能多地定向剩余的边。
\item \textbf{输出}：返回最终的 PAG。
\end{enumerate}

算法的计算复杂度主要取决于条件独立性检验的数量。对于 $p$ 个变量，最坏情况下需要检验的条件集数量是 $\sum_{k=0}^{p-2} \binom{p-2}{k}$，这使得 PC 算法在变量数很大时可能不切实际。

\section{其他因果发现方法概述}

PC 算法只是因果发现领域的众多种方法之一。这里简要介绍另外两类重要方法。

\subsection{基于评分的方法}

GIES（Greedy Equivalence类 Search）等方法将因果发现问题框架为模型选择问题。我们为每个候选 DAG 赋予一个评分（如 BIC 或 BD 评分），然后搜索评分最高的结构。

这类方法的优点是可以捕获全局最优结构；缺点是计算复杂度高，需要使用贪婪搜索等启发式方法。

\subsection{利用因果充分性假设}

许多方法假设不存在未观察到的共同原因（因果充分性）。在这个假设下，观测分布的每个因子对应于一个局部 DAG。

在这种假设下，因果方向可以通过检验的条件独立性完全确定。

\section{因果发现的局限性}

尽管因果发现方法取得了显著进展，我们必须清醒地认识其局限性。

\textbf{第一，假设的脆弱性。} 因果马尔可夫条件和忠诚性条件是强有力的假设。如果数据生成过程不满足这些条件（如存在未观察到的混杂），因果发现的结果可能是错误的。

\textbf{第二，检验的局限性。} 条件独立性检验的统计功效有限，特别是在小样本或大条件集时。可能出现假阳性（错误地删除边）或假阴性（错误地保留边）。

\textbf{第三，计算的复杂性。} 精确的因果发现是 NP 难的。现有算法都依赖于各种近似，这可能影响结果的可靠性。

\textbf{第四，不可辨识性。} 即使在理想条件下，某些因果结构也是不可辨识的——不同的因果图可能产生完全相同的观察分布。

\section{讨论与延伸阅读}

因果发现是一个活跃的研究领域，涉及统计学、计算机科学、哲学的交叉。经典参考文献包括 Pearl (2000) 的《Causality》，以及 Spirtes et al. (2000) 的《Causation, Prediction, and Search》。

近年的发展包括：
\begin{itemize}
\item 利用干预数据改进因果发现；
\item 将因果发现与机器学习结合；
\item 处理隐变量和选择偏差的更复杂方法。
\end{itemize}

在实际应用中，因果发现应该谨慎使用。其结果更适合作为假设生成，而不是因果声明的最终结论。

\section{本章小结}

本章介绍了从观察数据中发现因果结构的可能性的基本思想。

\begin{enumerate}
\item 在因果马尔可夫条件和忠诚性条件下，条件独立性蕴含因果结构信息。
\item PC 算法通过系统的条件独立性检验和 v-结构识别，从数据中学习因果骨架和方向。
\item 因果发现的结果以 PAG（部分有向无环图）的形式表示，反映了因果结构的等价类。
\item 因果发现有重要局限性，包括假设的脆弱性、统计检验的局限、计算复杂性和根本性的不可辨识性。
\end{enumerate}

下一步的学习方向包括：深入研究具体的条件独立性检验方法、了解基于评分的方法、以及探索利用干预数据的因果发现技术。

\section{习题}

\begin{Exercise}
证明：在马尔可夫条件和忠诚性条件下，如果 $X$ 和 $Y$ 在骨架中无边，但都与 $Z$ 相连，且 $X \not\Perp Y \mid \emptyset$ 但 $X \Perp Y \mid Z$，则 $(X, Z, Y)$ 形成 v-结构。
\end{Exercise}

\begin{Exercise}
考虑四个变量 $X_1, X_2, X_3, X_4$，假设真实因果结构是 $X_1 \to X_2 \to X_4$ 和 $X_1 \to X_3 \to X_4$。说明 PC 算法在什么条件下能够恢复这个结构，以及在什么条件下可能失败。
\end{Exercise}

\begin{Exercise}
证明：对于 $n$ 个变量的 DAG，完全图（有 $\binom{n}{2}$ 条边）和完全无向图在骨架发现步骤后的边数关系。
\end{Exercise}

\begin{Exercise}
讨论因果发现与因果推断其他方法（如第 10-12 章的倾向得分方法）的关系。它们分别适合什么场景？
\end{Exercise}
